\documentclass[11pt,a4paper]{article}
\usepackage[T1]{fontenc}
\usepackage[utf8]{inputenc}
\usepackage[brazil]{babel}
\usepackage{newpxtext,newpxmath}
\renewcommand\ttdefault{txtt}
\usepackage[a4paper,margin=25mm]{geometry}
\usepackage{amsmath}
\usepackage{microtype}
\usepackage[authoryear,round]{natbib}
\setcitestyle{aysep={},yysep={;}}
\usepackage{xurl}
\usepackage[hidelinks]{hyperref}
\hypersetup{pdftitle={Mudança na composição agrícola e cobertura da terra na Amazônia Legal (2000–2024)},pdfauthor={Rafael Eiki Teruya}}
\linespread{1.2}
\renewcommand\refname{Referências}

\begin{document}

\title{\bfseries Mudança na composição agrícola e cobertura da terra\\ na Amazônia Legal (2000–2024)}
\author{Rafael Eiki Teruya, N.º USP 14604204\\ {\small EAE1317, Economia do Meio Ambiente, FEA/USP}\\ {\small Prof. Danilo Igliori}}
\date{\small Entrega parcial}
\maketitle

\section{Introdução}

Publicado pelo Instituto Escolhas, o relatório \textit{Muita soja e pouca mandioca} examina os sistemas alimentares de 15 recortes metropolitanos da Amazônia Legal (incluindo a região imediata de Rio Branco, equivalente à metropolitana do Acre). O estudo reúne dados sobre população, produção agrícola, comercialização, preços, orçamento público e políticas de segurança alimentar, com atenção às condições de abastecimento nos centros urbanos. Na produção agrícola, compara a área plantada nas regiões entre 2017 e 2024 e registra movimentos em sentidos opostos: a área de soja cresceu 71\%, enquanto a de mandioca caiu 28\%. Em 2024, essas culturas correspondiam, respectivamente, a cerca de 62\% e 3\% da área plantada. O relatório também situa essa composição no contexto mais amplo das lavouras metropolitanas, destacando a concentração da área em poucas culturas e a presença de diferenças entre regiões e estados \citep{escolhas2026}. Partindo dessa motivação, este trabalho estende a análise para 2000–2024. Ele separa a retração absoluta da área alimentar da sua queda relativa, examina a distribuição municipal de perdas e ganhos e cruza essas mudanças com a cobertura da terra.

Por trás desses números há duas mudanças distintas: uma cultura perde participação quando sua área diminui, mas também quando as demais lavouras crescem mais depressa, ainda que a área dela permaneça estável ou aumente. Além disso, o saldo regional pode somar municípios em que a cultura recua e outros em que ela avança. Portanto, distinguir a retração absoluta da queda apenas relativa é o ponto de partida da análise.

Daí a pergunta central do projeto: \textit{como as mudanças na área agregada de nove culturas alimentares e na área conjunta de soja e milho contribuíram para a variação da participação alimentar nos municípios metropolitanos e nos demais municípios da Amazônia Legal, entre 2000 e 2024 e entre 2017 e 2024?} Nessa análise, mandioca, arroz, feijão, banana, abacaxi, batata-doce, melancia, laranja e mamão compõem o grupo de culturas alimentares considerado. A participação alimentar é a fração da área dessas nove culturas na área somada delas com soja e milho; ao contrário das participações do relatório, sua base não é a área plantada total. De forma complementar, examina-se a distribuição municipal de perdas e ganhos de área alimentar, as mudanças de vegetação nativa, pastagem e soja mapeada nos mesmos territórios e os custos de produção de mandioca e soja.

Trata-se de um estudo de caso empírico e descritivo, que combina a Produção Agrícola Municipal (PAM), do Instituto Brasileiro de Geografia e Estatística (IBGE), com as estatísticas de cobertura da terra do MapBiomas e os custos de produção da Conab. Para repartir a variação da participação alimentar entre a mudança da área dos alimentos e a de soja e milho, utiliza-se a decomposição de Shapley, baseada em um conceito da teoria dos jogos e sistematizada pelo economista Anthony Shorrocks para decompor indicadores econômicos entre seus componentes \citep{shorrocks2013}. O método é aplicado a duas janelas: a de 2017 a 2024, do relatório Escolhas, e a de 2000 a 2024, mais longa.

\section{Fundamentação econômica}\label{sec:fundamentacao}

Johann Heinrich von Thünen, proprietário rural no norte da Alemanha, formulou, na obra \textit{Der isolirte Staat in Beziehung auf Landwirthschaft und Nationalökonomie}, um modelo da localização das atividades agrícolas a partir da observação de sua própria fazenda \citep{thunen1826}. Ele supôs uma planície homogênea em torno de uma única cidade-mercado e concluiu que os usos da terra se organizam em anéis concêntricos, de acordo com o custo de levar cada produto até o mercado. Nessa formulação clássica, a escolha entre lavouras segue a renda que cada uso oferece à terra. Por hectare, a renda da cultura $i$ pode ser escrita como
\begin{equation}
R_i(d_i) = y_i\,[p_i - t_i d_i] - c_i,
\end{equation}
em que $y_i$ é a produtividade, $p_i$, o preço no mercado de destino da cultura $i$, $d_i$, a distância até esse mercado, $t_i$, o custo de transporte por unidade de produto e de distância, e $c_i$, os custos não fundiários por hectare. Este trabalho não estima $R_i$ para os municípios nem testa o efeito da distância aos mercados; o modelo serve para interpretar os resultados e orientar hipóteses. A comparação de custos da Seção~\ref{sec:metodo} apenas ilustra a ordem de grandeza da receita e de $c_i$ em dois sistemas de produção. Nesse quadro, a expansão de soja e milho é compatível com um aumento de sua renda relativa, seja por preço, produtividade, custos ou acesso a mercados. Como o mercado relevante muda com a cultura, uma hipótese para a leitura dos resultados é que a proximidade da capital pese mais para a mandioca, e portos e corredores de exportação, para os grãos. A classificação metropolitana delimita os grupos de comparação e não constitui uma medida direta de distância ou acesso aos mercados. Conta ainda a aglomeração: entre 1985 e 1995, a densidade populacional inicial de áreas comparáveis da Amazônia associou-se de forma não linear ao crescimento da produção agrícola e à conversão de terra. O efeito é positivo em níveis baixos de densidade e diminui até se inverter em níveis altos, por efeitos de congestão \citep{igliori2008}. Esse resultado motiva a análise separada dos municípios metropolitanos, uma escolha deste trabalho. No mesmo estudo, porém, a agricultura cresceu mais nas áreas com maior custo de transporte até a capital estadual mais próxima, o que relativiza o peso dos mercados urbanos locais na hipótese sobre a mandioca.

Quando emissões, perda de biodiversidade e efeitos hídricos recaem sobre terceiros sem preço, a decisão privada de converter vegetação deixa de incorporá-los \citep{baumol1988,kolstad2011}. Em uma representação simplificada, sendo $\Delta R$ o ganho privado de converter uma parcela e $\Delta D$ o dano externo incremental, a conversão interessa ao produtor sempre que $\Delta R > 0$, mas só é socialmente desejável se $\Delta R > \Delta D$. A expansão agrícola pode, ainda, deslocar pastagens e desmatamento para outros municípios \citep{arima2011,richards2014}. Os dados do MapBiomas permitem dimensionar as mudanças líquidas de cobertura em cada município, mas não medem $\Delta D$, não atribuem a conversão a uma cultura específica nem captam deslocamentos para outros municípios.

A conversão de terra no período também respondeu a regras públicas e a compromissos de mercado. A Moratória da Soja, firmada em 2006, comprometeu os compradores a não adquirir soja de áreas desmatadas após julho daquele ano \citep{nepstad2014}; a data de corte foi depois transferida para julho de 2008. A Resolução 3.545 do Conselho Monetário Nacional (2008) condicionou o crédito rural no bioma Amazônia à comprovação de regularidade ambiental \citep{bcb2008}, e a Lei 12.651 (2012) alterou as regras de regularização das áreas rurais \citep{brancalion2016}. Essas regras afetaram a expansão da soja e o desmatamento em datas e intensidades diferentes. A janela de 2000 a 2024 atravessa todas essas mudanças, enquanto a de 2017 a 2024 começa depois de todas elas. Sem estimar seus efeitos, o trabalho as trata como contexto das mudanças de área e de cobertura.

\section{Dados}\label{sec:dados}

O universo de análise reúne os 772 municípios da lista da Amazônia Legal de 2022 \citep{ibge2022}. O relatório Escolhas parte da lista de 2024, que também soma 772 municípios, e dele vem a classificação dos 138 municípios metropolitanos. Da tabela 5457 do Sistema IBGE de Recuperação Automática (SIDRA), base de consulta às estatísticas do IBGE, foram extraídos dados da Produção Agrícola Municipal sobre área plantada ou destinada à colheita, quantidade e valor da produção entre 2000 e 2024 \citep{pam2026}. Em 128 dos 19.300 municípios-ano, estão indisponíveis tanto a área total quanto as áreas das onze culturas consideradas (as nove alimentares, a soja e o milho), sendo que 36 desses casos são anteriores à instalação dos respectivos municípios. Esses registros são mantidos como ausentes, pois a indisponibilidade de informação não indica ausência de cultivo. Nos anos anteriores à instalação, o território ainda integrava outro município, de modo que atribuir zero poderia criar variações artificiais na área agrícola.

Para a análise principal, utilizam-se os 746 municípios observados nos 25 anos, dos quais 133 são metropolitanos e 613 não. Embora o método use apenas 2000, 2017 e 2024, a exigência da série completa restringe a análise a municípios com registros contínuos, já que lacunas na série podem indicar falhas na coleta da informação. Essa escolha exige excluir 26 municípios, cinco deles metropolitanos. Como teste de robustez, a decomposição é repetida com os 754 municípios observados nos três anos usados pelo método, o que recupera oito dos 26 excluídos: três metropolitanos e cinco não metropolitanos. Nas duas amostras, os municípios são os mesmos em todos os anos comparados. Como os municípios de origem permanecem na amostra, a instalação de um novo município reduz artificialmente a área agrícola do município de origem a partir desse ano; esse efeito é reconhecido como limitação da análise.

Define-se $F$ como a área somada das nove culturas selecionadas e $G$ como a área de soja e milho, de modo que a participação alimentar é $s = F/(F+G)$. O grupo $F$ reúne grãos (arroz e feijão), raízes (mandioca e batata-doce) e frutas (banana, abacaxi, melancia, laranja e mamão), em um recorte operacional para acompanhar mudanças na composição agrícola. Todas essas culturas aparecem na análise do Instituto Escolhas, mas sua reunião em um único indicador é uma escolha deste trabalho, não uma classificação adotada pelo relatório. Nas regiões metropolitanas, oito das nove culturas perderam área entre 2017 e 2024, e mandioca e arroz respondiam, em 2024, por cerca de 70\% da área do grupo. Já as culturas alimentares excluídas que cresceram, como tomate e limão, somariam menos de 1\% da área do grupo \citep{escolhas2026}; por isso, incluí-las tenderia a alterar pouco $s$. Esse argumento se apoia nos dados metropolitanos de 2017 a 2024; para os demais municípios e para a janela de 2000 a 2024, a área dessas culturas é verificada na PAM. O grupo não pretende representar toda a produção alimentar ou a dieta regional. O açaí fica de fora por não ter série completa na PAM desde 2000. Como as áreas são somadas dentro de cada grupo de municípios antes do cálculo, municípios com mais área pesam mais. Em um segundo teste, $G$ é substituído por $O = \text{Total PAM} - F$, que inclui todas as outras lavouras. A PAM informa a área plantada ou destinada à colheita. Para mandioca e abacaxi, culturas de longa duração, considera-se a área destinada à colheita no ano. Como uma mesma parcela pode receber cultivos sucessivos ou simultâneos, a soma das áreas das culturas não corresponde necessariamente à superfície física ocupada. Isso afeta sobretudo $G$, já que o milho de segunda safra ocupa em grande parte a mesma área da soja, de modo que o aumento de $G$ não implica, necessariamente, ocupação de terra nova.

Para a cobertura da terra, utiliza-se a Coleção 10.1 do MapBiomas, uma iniciativa colaborativa que produz mapas anuais de cobertura e uso da terra no Brasil a partir de imagens de satélite, classificadas por métodos computacionais \citep{mapbiomas2026}. Vegetação nativa reúne as formações florestais (classes 3, 4, 5, 6 e 49) e não florestais (classes 11, 12, 29, 32 e 50) da legenda da Coleção 10, incluindo a vegetação secundária quando classificada nessas formações, mas não necessariamente toda a área em regeneração. Pastagem e soja correspondem às classes 15 e 39; a descrição da classe soja refere-se à monocultura de primeira safra. Foram mantidos os mesmos 746 códigos e os territórios municipais inteiros, inclusive as partes de 21 municípios maranhenses situadas fora do polígono legal. A soja mapeada não equivale a $G$, que inclui milho, e sua área não é diretamente equivalente à informada pela PAM para essa cultura. O total mapeado por município é praticamente constante no período, mas isso não garante que os limites coincidam com os usados pela PAM em cada ano. Como as duas fontes não são ajustadas a territórios comuns ao longo do tempo, eventuais diferenças de limites permanecem como uma limitação cujo efeito não é quantificado.

Já os custos vêm das planilhas da Companhia Nacional de Abastecimento (Conab), empresa pública federal que acompanha a produção e o abastecimento agropecuário. Essas planilhas estimam os custos por hectare para sistemas de produção definidos em localidades específicas, com base nas quantidades de insumos e serviços utilizados e em seus preços \citep{conab2020}. Utilizam-se as planilhas de março de 2024 para a safra 2024/25 \citep{conab2024}, com dois orçamentos do Maranhão: mandioca de agricultura familiar e baixa tecnologia em Santa Luzia e soja empresarial de alta tecnologia, em plantio direto, em Balsas. Como os dois orçamentos diferem também em escala e tecnologia, a comparação não isola o efeito da cultura. A receita vem da PAM de 2024 nos mesmos municípios. Como a PAM registra a colheita do ano civil, a soja de 2024 corresponde à safra 2023/24, e não à 2024/25 dos orçamentos; a comparação é, portanto, a preços de 2024, e não de uma mesma safra.

\section{Método}\label{sec:metodo}

No artigo \textit{Decomposition procedures for distributional analysis: a unified framework based on the Shapley value}, \citet{shorrocks2013} apresenta uma estrutura comum para decompor indicadores na análise distributiva, como medidas de desigualdade entre diferentes fontes de renda. O problema consiste em repartir um indicador entre seus componentes quando a contribuição de cada componente depende dos demais. Na teoria dos jogos cooperativos, o valor de Shapley reparte o resultado da cooperação entre os participantes, atribuindo a cada um a média de sua contribuição marginal em todas as ordens possíveis de entrada no grupo. Shorrocks adapta essa regra à decomposição de indicadores, em que a contribuição de cada componente é calculada pela média de seus efeitos marginais em todas as ordens possíveis de inclusão. Assim, a atribuição não depende de uma ordem específica escolhida pelo pesquisador.

A aplicação considera a variação temporal da participação alimentar, definida como $h(F,G) = F/(F+G)$, em que $F$ é a área das nove culturas selecionadas e $G$, a área conjunta de soja e milho. Como a contribuição de cada área depende da outra, a ordem de atualização altera as parcelas atribuídas a $F$ e $G$. Calculam-se, por isso, os efeitos marginais nas duas ordens possíveis e toma-se sua média, conforme as expressões abaixo. As contribuições somam a variação observada, sem que a decomposição implique uma sequência histórica ou uma relação causal. Com subscritos 0 (inicial) e 1 (final),
\begin{align}
C_F &= \tfrac{1}{2}\,[h(F_1,G_0) - h(F_0,G_0) + h(F_1,G_1) - h(F_0,G_1)]\,, \\
C_G &= \tfrac{1}{2}\,[h(F_0,G_1) - h(F_0,G_0) + h(F_1,G_1) - h(F_1,G_0)]\,,
\end{align}
de modo que $C_F + C_G = h(F_1,G_1) - h(F_0,G_0) = \Delta s$. Para ilustrar, se $F$ cai de 100 para 50 hectares e $G$ sobe de 100 para 200, a participação cai 30 pontos percentuais; dependendo da ordem, a contribuição dos alimentos é de $-16{,}7$ ou $-13{,}3$ pontos, e Shapley atribui $-15$ a cada componente.

Aplica-se a decomposição às áreas agregadas de cada grupo e janela, pois, como $h$ é não linear, somar decomposições municipais levaria a outro resultado. Além das contribuições em pontos percentuais, apresentam-se a faixa entre as duas ordens e, quando $C_F$ e $C_G$ têm o mesmo sinal, a parcela $C_G/\Delta s$, que perde sentido quando os sinais diferem ou quando $\Delta s$ é próximo de zero. Essa faixa mede a sensibilidade à regra de atribuição e não deve ser lida como intervalo de confiança. Os estados intermediários, como $(F_1, G_0)$, são apenas combinações aritméticas, não cenários que um produtor poderia escolher. Pela dupla contagem descrita na Seção~\ref{sec:dados}, parte de $C_G$ pode refletir o milho de segunda safra plantado na área da soja, e não espaço tomado dos alimentos; o teste com $O$ não elimina esse efeito, pois $O$ também inclui o milho.

Quanto à distribuição municipal, os municípios são classificados, em cada janela, pelo sinal de $\Delta F_m = F_{m1} - F_{m0}$ (os casos com $\Delta F_m = 0$ são contados à parte), e as perdas daqueles em retração são somadas separadamente dos ganhos daqueles em expansão. Comparam-se, então, vegetação, pastagem e soja mapeada entre esses grupos, calculando a variação proporcional sobre os estoques agregados. Como são definidos depois de observar $\Delta F_m$, os grupos servem para descrever coexistência territorial, não para estimar efeitos.

No caso dos custos, aproxima-se $c_i$ da equação (1) pelo custo total da Conab, descontando a remuneração da terra (terra própria e arrendamento). Toma-se como $p_i - t_i d_i$ da equação (1) o preço implícito da PAM, obtido pela razão entre valor e quantidade produzidos, que é o preço recebido pelo produtor. A partir dele, a receita por hectare é calculada sob duas produtividades, a da PAM e a do orçamento, e obtém-se também a produtividade em que a receita iguala $c_i$.

\section{Próximas etapas}

Para a entrega final, a próxima etapa é aplicar aos dados o método descrito na Seção~\ref{sec:metodo}. Primeiro, calculam-se $F$, $G$ e $s$ para os municípios metropolitanos e para os demais, em 2000, 2017 e 2024, junto com a área das culturas alimentares excluídas de $F$. Em seguida, aplica-se a decomposição de Shapley às duas janelas, com $C_F$, $C_G$, a faixa entre as duas ordens e, quando os sinais coincidirem, a parcela $C_G/\Delta s$. A decomposição é então repetida nos dois testes já definidos: com os 754 municípios observados em 2000, 2017 e 2024, e com $O = \text{Total PAM} - F$ no lugar de $G$. Na distribuição municipal, os municípios são classificados, em cada janela, pelo sinal de $\Delta F_m$. As perdas e os ganhos são somados separadamente, e comparam-se entre os grupos as variações proporcionais de vegetação nativa, pastagem e soja mapeada do MapBiomas. Por fim, com os orçamentos da Conab de mandioca em Santa Luzia e de soja em Balsas e com os preços implícitos da PAM de 2024, calculam-se a receita por hectare sob as duas produtividades e a produtividade em que a receita iguala $c_i$. Os resultados serão lidos à luz da Seção~\ref{sec:fundamentacao}, como descrição de mudanças de composição e de cobertura, sem atribuição causal.

\section*{Uso de inteligência artificial}

Utilizou-se a inteligência artificial como apoio à elaboração e à revisão de scripts para o tratamento dos dados e à localização de tabelas e documentos das fontes consultadas; na entrega final, ela também apoiará os cálculos. Na escrita, auxiliou na revisão da clareza, da organização e da correção gramatical do texto. Os resultados fornecidos pela IA foram conferidos com os dados e com a documentação oficial disponível, e seu uso limitou-se a esse apoio.

\begin{thebibliography}{99}

\bibitem[Arima et~al.(2011)Arima, Richards, Walker e Caldas]{arima2011}
Arima, E.~Y., Richards, P., Walker, R. e Caldas, M.~M. (2011).
\newblock Statistical confirmation of indirect land use change in the Brazilian Amazon.
\newblock \textit{Environmental Research Letters}, 6(2):024010.

\bibitem[{Banco Central do Brasil}(2008)]{bcb2008}
Banco Central do Brasil (2008).
\newblock Resolução n.º 3.545, de 29 de fevereiro de 2008.
\newblock Conselho Monetário Nacional. Disponível em: \url{https://www.bcb.gov.br/pre/normativos/res/2008/pdf/res_3545_v1_O.pdf}. Acesso em: 2 out. 2026.

\bibitem[Baumol e Oates(1988)]{baumol1988}
Baumol, W.~J. e Oates, W.~E. (1988).
\newblock \textit{The Theory of Environmental Policy}.
\newblock 2. ed. Cambridge University Press, Cambridge.

\bibitem[Brancalion et~al.(2016)Brancalion, Garcia, Loyola, Rodrigues, Pillar e Lewinsohn]{brancalion2016}
Brancalion, P. H.~S., Garcia, L.~C., Loyola, R., Rodrigues, R.~R., Pillar, V.~D. e Lewinsohn, T.~M. (2016).
\newblock A critical analysis of the Native Vegetation Protection Law of Brazil (2012): updates and ongoing initiatives.
\newblock \textit{Natureza \& Conservação}, 14:1--15.

\bibitem[{Companhia Nacional de Abastecimento}(2020)]{conab2020}
Companhia Nacional de Abastecimento (2020).
\newblock Norma 30.302: metodologia de custo de produção.
\newblock Disponível em: \url{https://www.gov.br/conab/pt-br/acesso-a-informacao/institucional/atos-normativos/normas-da-organizacao/operacoes/30-302_norma_metodologia_de_custo_de_producao.pdf}. Acesso em: 2 out. 2026.

\bibitem[{Companhia Nacional de Abastecimento}(2024)]{conab2024}
Companhia Nacional de Abastecimento (2024).
\newblock Custos de produção agrícolas: séries históricas de mandioca e soja.
\newblock Abas Santa Luzia-MA-2024 e Balsas-MA-2024, preços de março de 2024, safra 2024/25. Disponível em: \url{https://www.gov.br/conab/pt-br/atuacao/informacoes-agropecuarias/custos-de-producao/planilhas-de-custos-de-producao}. Acesso em: 2 out. 2026.

\bibitem[Igliori(2008)]{igliori2008}
Igliori, D.~C. (2008).
\newblock Deforestation, growth and agglomeration effects: Evidence from agriculture in the Brazilian Amazon.
\newblock Environmental Economy and Policy Research Discussion Paper 28.2008, Department of Land Economy, University of Cambridge.

\bibitem[{Instituto Brasileiro de Geografia e Estatística}(2022)]{ibge2022}
Instituto Brasileiro de Geografia e Estatística (2022).
\newblock Municípios da Amazônia Legal 2022.
\newblock Lista territorial oficial. Disponível em: \url{https://geoftp.ibge.gov.br/organizacao_do_territorio/estrutura_territorial/amazonia_legal/2022/Municipios_da_Amazonia_Legal_2022.ods}. Acesso em: 2 out. 2026.

\bibitem[{Instituto Brasileiro de Geografia e Estatística}(2026)]{pam2026}
Instituto Brasileiro de Geografia e Estatística (2026).
\newblock Produção agrícola municipal: tabela 5457 -- área plantada, quantidade produzida e valor da produção das lavouras temporárias e permanentes, 2000--2024.
\newblock Sistema IBGE de Recuperação Automática (SIDRA). Disponível em: \url{https://sidra.ibge.gov.br/tabela/5457}. Acesso em: 2 out. 2026.

\bibitem[{Instituto Escolhas}(2026)]{escolhas2026}
Instituto Escolhas (2026).
\newblock Muita soja e pouca mandioca: a alimentação nas regiões metropolitanas da Amazônia Legal.
\newblock Relatório técnico, Instituto Escolhas, São Paulo. Publicado em set. 2026. Disponível em: \url{https://escolhas.org/publicacao/muita-soja-e-pouca-mandioca-a-alimentacao-nas-regioes-metropolitanas-da-amazonia-legal/}. Acesso em: 2 out. 2026.

\bibitem[Kolstad(2011)]{kolstad2011}
Kolstad, C.~D. (2011).
\newblock \textit{Environmental Economics}.
\newblock 2. ed. Oxford University Press, New York.

\bibitem[Nepstad et~al.(2014)]{nepstad2014}
Nepstad, D., McGrath, D., Stickler, C., Alencar, A., Azevedo, A., Swette, B., Bezerra, T., DiGiano, M., Shimada, J., Seroa da Motta, R., Armijo, E., Castello, L., Brando, P., Hansen, M.~C., McGrath-Horn, M., Carvalho, O. e Hess, L. (2014).
\newblock Slowing Amazon deforestation through public policy and interventions in beef and soy supply chains.
\newblock \textit{Science}, 344(6188):1118--1123.

\bibitem[{Projeto MapBiomas}(2026)]{mapbiomas2026}
Projeto MapBiomas (2026).
\newblock Estatísticas de cobertura e uso da terra por município, estado e bioma: Coleção 10.1, 1985--2024.
\newblock Versão 10.1, publicada em 19 fev. 2026. Disponível em: \url{https://doi.org/10.58053/MapBiomas/SJZOLT}. Acesso em: 2 out. 2026.

\bibitem[Richards et~al.(2014)Richards, Walker e Arima]{richards2014}
Richards, P.~D., Walker, R.~T. e Arima, E.~Y. (2014).
\newblock Spatially complex land change: The indirect effect of Brazil's agricultural sector on land use in Amazonia.
\newblock \textit{Global Environmental Change}, 29:1--9.

\bibitem[Shorrocks(2013)]{shorrocks2013}
Shorrocks, A.~F. (2013).
\newblock Decomposition procedures for distributional analysis: a unified framework based on the Shapley value.
\newblock \textit{The Journal of Economic Inequality}, 11(1):99--126.

\bibitem[von Thünen(1826)]{thunen1826}
von Thünen, J.~H. (1826).
\newblock \textit{Der isolirte Staat in Beziehung auf Landwirthschaft und Nationalökonomie}.
\newblock Perthes, Hamburg.

\end{thebibliography}

\end{document}
